This work studies inference efficiency for existing language models. It does not train or release a new target model, does not change any target's output distribution under greedy verification, and does not introduce new capabilities or new training data about people. All targets, drafters and datasets are publicly released artifacts used under their own licenses, and the domain adapters we evaluate come from previously released checkpoints.

The main societal consideration is indirect. Faster inference lowers the cost of deploying language models, which amplifies both their benefits and any harms of the underlying target model. RelaySpec inherits whatever behavior, bias or failure mode its target already has, and speculative decoding provides no additional safety property. Because adaptation is cheap, it also lowers the cost of maintaining many specialized target variants, which makes the provenance of each adapter worth tracking in a deployment. We report negative and partial results, including the workloads and model pairs where the method does not reach a native drafter, so that the method's limits are visible to anyone considering it.
